\begin{thebibliography}{36}
\bibitem{ref1}Miklós Ajtai. Σ11 -formulae on finite structures. Ann. Pure Appl. Logic, 24(1):1–48, 1983. doi:10.1016/0168-0072(83)90038-6.
\bibitem{ref2}Adi Akavia, Andrej Bogdanov, Siyao Guo, Akshay Kamath, and Alon Rosen. Candidate weak pseudorandom functions in AC0 ◦ MOD2 . In Innovations in Theoretical Computer Science, ITCS’14, Princeton, NJ, USA, January 12-14, 2014, pages 251–260, 2014. doi: 10.1145/2554797.2554821.
\bibitem{ref3}Josh Alman and R. Ryan Williams. Probabilistic rank and matrix rigidity. In Proceedings of the 49th Annual ACM SIGACT Symposium on Theory of Computing, STOC 2017, Montreal, QC, Canada, June 19-23, 2017, pages 641–652, 2017. doi:10.1145/3055399.3055484.
\bibitem{ref4}László Babai, Peter Frankl, and Janos Simon. Complexity classes in communication complexity theory (preliminary version). In 27th Annual Symposium on Foundations of Computer Science, Toronto, Canada, 27-29 October 1986, pages 337–347, 1986. doi:10.1109/SFCS.1986.15.
\bibitem{ref5}László Babai, Anna Gál, Peter G. Kimmel, and Satyanarayana V. Lokam. Communication complexity of simultaneous messages. SIAM J. Comput., 33(1):137–166, 2003. doi:10.1137/ S0097539700375944.
\bibitem{ref6}Abhishek Banerjee, Chris Peikert, and Alon Rosen. Pseudorandom functions and lattices. In Advances in Cryptology - EUROCRYPT 2012 - 31st Annual International Conference on the Theory and Applications of Cryptographic Techniques, Cambridge, UK, April 15-19, 2012. Proceedings, pages 719–737, 2012. doi:10.1007/978-3-642-29011-4\_42.
\bibitem{ref7}Avrim Blum, Merrick L. Furst, Jeffrey C. Jackson, Michael J. Kearns, Yishay Mansour, and Steven Rudich. Weakly learning DNF and characterizing statistical query learning using fourier analysis. In Proceedings of the Twenty-Sixth Annual ACM Symposium on Theory of Computing, 23-25 May 1994, Montréal, Québec, Canada, pages 253–262. ACM, 1994. doi:10.1145/195058.195147.
\bibitem{ref8}Andrej Bogdanov, Yuval Ishai, and Akshayaram Srinivasan. Unconditionally secure compu- tation against low-complexity leakage. In Advances in Cryptology - CRYPTO 2019 - 39th Annual International Cryptology Conference, Santa Barbara, CA, USA, August 18-22, 2019, Proceedings, Part II, volume 11693 of Lecture Notes in Computer Science, pages 387–416, 2019. doi:10.1007/978-3-030-26951-7\_14.
\bibitem{ref9}Andrej Bogdanov and Alon Rosen. Pseudorandom functions: Three decades later. In Tutorials on the Foundations of Cryptography, pages 79–158. Springer, Cham, 2017. doi: 10.1007/978-3-319-57048-8\_3.
\bibitem{ref10}Dan Boneh, Yuval Ishai, Alain Passelègue, Amit Sahai, and David J. Wu. Exploring crypto dark matter: - new simple PRF candidates and their applications. In Theory of Cryptography - 16th International Conference, TCC 2018, Panaji, India, November 11-14, 2018, Proceedings, Part II, volume 11240 of Lecture Notes in Computer Science, pages 699–729, 2018. doi: 10.1007/978-3-030-03810-6\_25.
\bibitem{ref11}Arkadev Chattopadhyay and Rahul Santhanam. Lower bounds on interactive compressibility by constant-depth circuits. In 53rd Annual IEEE Symposium on Foundations of Computer Science, FOCS 2012, New Brunswick, NJ, USA, October 20-23, 2012, pages 619–628, 2012. doi:10.1109/FOCS.2012.74.
\bibitem{ref12}Mahdi Cheraghchi, Elena Grigorescu, Brendan Juba, Karl Wimmer, and Ning Xie. AC0 ◦MOD2 lower bounds for the boolean inner product. In 43rd International Colloquium on Automata, Languages, and Programming, ICALP 2016, July 11-15, 2016, Rome, Italy, volume 55 of LIPIcs, pages 35:1–35:14, 2016. doi:10.4230/LIPIcs.ICALP.2016.35.
\bibitem{ref13}Amit Daniely, Nati Linial, and Shai Shalev-Shwartz. From average case complexity to improper learning complexity. In Symposium on Theory of Computing, STOC 2014, New York, NY, USA, May 31 - June 03, 2014, pages 441–448, 2014. doi:10.1145/2591796.2591820.
\bibitem{ref14}Ning Ding, Yanli Ren, and Dawu Gu. PAC learning depth-3 AC0 circuits of bounded top fanin. In International Conference on Algorithmic Learning Theory, ALT 2017, 15-17 October 2017, Kyoto University, Kyoto, Japan, pages 667–680, 2017. URL: http://proceedings.mlr. press/v76/ding17a.html.
\bibitem{ref15}Bella Dubrov and Yuval Ishai. On the randomness complexity of efficient sampling. In Proceedings of the 38th Annual ACM Symposium on Theory of Computing, Seattle, WA, USA, May 21-23, 2006, pages 711–720, 2006. doi:10.1145/1132516.1132615.
\bibitem{ref16}Zeev Dvir and Benjamin Edelman. Matrix rigidity and the croot-lev-pach lemma. CoRR, abs/1708.01646, 2017. arXiv:1708.01646.
\bibitem{ref17}Merrick L. Furst, James B. Saxe, and Michael Sipser. Parity, circuits, and the polynomial- time hierarchy. In 22nd Annual Symposium on Foundations of Computer Science, Nashville, Tennessee, USA, 28-30 October 1981, pages 260–270, 1981. doi:10.1109/SFCS.1981.35.
\bibitem{ref18}Edgar N. Gilbert. A comparison of signalling alphabets. The Bell system technical journal, 31(3):504–522, 1952.
\bibitem{ref19}Oded Goldreich, Shafi Goldwasser, and Silvio Micali. How to construct random functions. J. ACM, 33(4):792–807, 1986. doi:10.1145/6490.6503.
\bibitem{ref20}Mika Göös, Toniann Pitassi, and Thomas Watson. The landscape of communication complexity classes. In 43rd International Colloquium on Automata, Languages, and Programming, ICALP 2016, July 11-15, 2016, Rome, Italy, volume 55 of LIPIcs, pages 86:1–86:15, 2016. doi: 10.4230/LIPIcs.ICALP.2016.86.
\bibitem{ref21}Robert M. Gray. Entropy and information theory. Springer Science \& Business Media, 2011.
\bibitem{ref22}Johan Håstad. Almost optimal lower bounds for small depth circuits. In Proceedings of the 18th Annual ACM Symposium on Theory of Computing, May 28-30, 1986, Berkeley, California, USA, pages 6–20, 1986. doi:10.1145/12130.12132.
\bibitem{ref23}Michael J. Kearns. Efficient noise-tolerant learning from statistical queries. In Proceedings of the Twenty-Fifth Annual ACM Symposium on Theory of Computing, May 16-18, 1993, San Diego, CA, USA, pages 392–401. ACM, 1993. doi:10.1145/167088.167200.
\bibitem{ref24}Nathan Linial, Yishay Mansour, and Noam Nisan. Constant depth circuits, fourier transform, and learnability. In 30th Annual Symposium on Foundations of Computer Science, Research Triangle Park, North Carolina, USA, 30 October - 1 November 1989, pages 574–579, 1989. doi:10.1109/SFCS.1989.63537.
\bibitem{ref25}Pavel Pudlák, Vojtech Rödl, and Petr Savický. Graph complexity. Acta Inf., 25(5):515–535, 1988. doi:10.1007/BF00279952.
\bibitem{ref26}Alexander A. Razborov. On rigid matrices. preprint, 1989.
\bibitem{ref27}Guy N. Rothblum. How to compute under AC0 leakage without secure hardware. In Advances in Cryptology - CRYPTO 2012 - 32nd Annual Cryptology Conference, Santa Barbara, CA, USA, August 19-23, 2012. Proceedings, volume 7417 of Lecture Notes in Computer Science, pages 552–569, 2012. doi:10.1007/978-3-642-32009-5\_32.
\bibitem{ref28}Rocco A. Servedio and Emanuele Viola. On a special case of rigidity. Electronic Colloquium on Computational Complexity (ECCC), 19:144, 2012. URL: http://eccc.hpi-web.de/report/ 2012/144.
\bibitem{ref29}Avishay Tal. The bipartite formula complexity of inner-product is quadratic. Electronic Colloquium on Computational Complexity (ECCC), 23:181, 2016. URL: http://eccc.hpi-web. de/report/2016/181.
\bibitem{ref30}Avishay Tal. Tight bounds on the fourier spectrum of AC0 . In 32nd Computational Complexity Conference, CCC 2017, July 6-9, 2017, Riga, Latvia, volume 79 of LIPIcs, pages 15:1–15:31, 2017. doi:10.4230/LIPIcs.CCC.2017.15.
\bibitem{ref31}Salil P. Vadhan. On learning vs. refutation. In Proceedings of the 30th Conference on Learning Theory, COLT 2017, Amsterdam, The Netherlands, 7-10 July 2017, pages 1835–1848, 2017. URL: http://proceedings.mlr.press/v65/vadhan17a.html.
\bibitem{ref32}Leslie G. Valiant. A theory of the learnable. In Proceedings of the 16th Annual ACM Symposium on Theory of Computing, April 30 - May 2, 1984, Washington, DC, USA, pages 436–445, 1984. doi:10.1145/800057.808710.
\bibitem{ref33}R. R. Varshamov. Estimate of the number of signals in error correcting codes. Docklady Akad. Nauk, SSSR, 117:739–741, 1957.
\bibitem{ref34}Emanuele Viola. Extractors for circuit sources. In IEEE 52nd Annual Symposium on Founda- tions of Computer Science, FOCS 2011, Palm Springs, CA, USA, October 22-25, 2011, pages 220–229, 2011. doi:10.1109/FOCS.2011.20.
\bibitem{ref35}Henning Wunderlich. On a theorem of razborov. Computational Complexity, 21(3):431–477, 2012. doi:10.1007/s00037-011-0021-5.
\bibitem{ref36}Andrew Chi-Chih Yao. Some complexity questions related to distributive computing (prelimin- ary report). In Proceedings of the 11h Annual ACM Symposium on Theory of Computing, April 30 - May 2, 1979, Atlanta, Georgia, USA, pages 209–213, 1979. doi:10.1145/800135.804414.
\end{thebibliography}
